\documentclass[11pt,a4paper]{article}
\usepackage[T1]{fontenc}
\usepackage[utf8]{inputenc}
\usepackage{lmodern,microtype}
\usepackage[a4paper,margin=25mm,headheight=15pt]{geometry}
\usepackage{amsmath,amssymb,amsthm,mathtools,mathrsfs}
\usepackage{booktabs,longtable,array,enumitem}
\usepackage[dvipsnames]{xcolor}
\usepackage{fancyhdr,titlesec}
\usepackage[unicode,colorlinks=true,linkcolor=MidnightBlue,citecolor=MidnightBlue,urlcolor=MidnightBlue,breaklinks=true]{hyperref}
\usepackage{bookmark}
\usepackage{listings}
\definecolor{ink}{RGB}{25,47,65}
\titleformat{\section}{\large\bfseries\color{ink}}{\thesection}{0.65em}{}
\titleformat{\subsection}{\normalsize\bfseries\color{ink}}{\thesubsection}{0.65em}{}
\pagestyle{fancy}\fancyhf{}
\fancyhead[L]{\small Negative-ray summation of feedback transseries}
\fancyhead[R]{\small Research article}
\fancyfoot[C]{\thepage}
\renewcommand{\headrulewidth}{0.3pt}
\setlength{\parskip}{3pt}
\setlength{\emergencystretch}{3em}
\setlist{itemsep=4pt,topsep=5pt,leftmargin=1.7em}
\numberwithin{equation}{section}
\newtheorem{theorem}{Theorem}[section]
\newtheorem{proposition}[theorem]{Proposition}
\newtheorem{lemma}[theorem]{Lemma}
\newtheorem{corollary}[theorem]{Corollary}
\theoremstyle{definition}
\newtheorem{definition}[theorem]{Definition}
\newtheorem{example}[theorem]{Example}
\theoremstyle{remark}
\newtheorem{remark}[theorem]{Remark}
\newcommand{\C}{\mathbb C}
\newcommand{\R}{\mathbb R}
\newcommand{\N}{\mathbb N}
\newcommand{\U}{\widehat U}
\newcommand{\Q}{\widehat Q}
\newcommand{\Pp}{\widehat P}
\newcommand{\B}{\mathcal B_0}
\newcommand{\G}{\mathcal G_1}
\newcommand{\ee}{\mathrm e}
\newcommand{\dd}{\mathrm d}
\newcommand{\Ree}{\operatorname{Re}}
\newcommand{\Imm}{\operatorname{Im}}
\newcommand{\dist}{\operatorname{dist}}
\newcommand{\supp}{\operatorname{supp}}
\newcommand{\inv}{\langle-1\rangle}
\newcommand{\TV}{\mathrm{TV}}
\newcommand{\coeff}[2]{[#1^{#2}]}
\newcommand{\sh}{\mathsf C}
\hypersetup{pdftitle={Negative-Ray Summation of Countable Exponential-Feedback Transseries},pdfauthor={Research draft prepared with ChatGPT for Vladimir Reshetnikov},pdfsubject={Fine Borel summability, inverse Bessel kernels, logarithmic-Gevrey remainders, certified inversion}}
% ed. (2026-09-29): editorial-note environment for ProveIt cross-references;
% it uses no counter, so no author numbering changes.
\newenvironment{ednote}{\par\smallskip\begin{quote}\small\raggedright\textbf{Editorial note (ProveIt, 2026-09-29).}\ }{\end{quote}\smallskip}
\begin{document}
\begin{center}
{\LARGE\bfseries\color{ink} Negative-Ray Summation of Countable\par Exponential-Feedback Transseries\par}
\vspace{0.7em}
{\large A sharp slope threshold, an inverse Bessel representation,\par and logarithmic-Gevrey error certificates\par}
\vspace{0.9em}
Research article prepared for Vladimir Reshetnikov\\
29 September 2026
\end{center}

\begin{abstract}
Consider the formally well-defined feedback equation
\[
 \U(q)=\sum_{j\ge1}q^j\exp(\lambda_j\U(q)),\qquad \lambda_j\ge0,
\]
and its compositional inverse $\Q$.  We prove that $\U$ and $\Q$ are
\emph{finely Borel--Laplace summable in the negative direction} if and only
if $\lambda_j=O((j\log(j+1))^2)$.  The sums are inverse functions and recover
the literal feedback equation near the negative real axis.  The proof
inverts the feedback equation before taking its Borel transform: an exact
signed exponential-block expansion becomes a normally convergent sum of
Bessel kernels.  If
$A=\limsup\lambda_j/(j\log j)^2<\infty$, this constructs the inverse Borel
transform on the parabolic region
$|\zeta|+\Ree\zeta<1/(2A)$; for $A=0$ it constructs an entire function.
For the quadratic model $\lambda_j=j^2$, we additionally prove a uniform
remainder estimate on a closed left half-disc.  Its explicit positive
majorant $M_N$ satisfies
\[
 \log M_N=N\log N-2N\log\log N+(\log4-1)N+o(N).
\]
Truncation at $N\sim\log^2(1/r)/(4r)$ consequently gives an error at most
$\exp[-(1/4-o(1))\log^2(1/r)/r]$.  We provide exact rational numerical
certificates, an inverse-error transport bound, reproducible programs,
and a research agenda separating fine summability from angular summability,
resurgence, and actual optimal-error asymptotics.
\end{abstract}

\noindent\textbf{Status.} The theorems below are conventional mathematical proofs,
not Lean-verified statements.  The proposed new contribution is the constructive
negative-direction summation theorem and its quantitative remainder theory for
this model.  The underlying Lagrange, Bessel, and nonlinear summation tools are
classical.  Global publication priority has not been established.  In particular,
we do not claim a general transseries summation theorem or a solution of a named
longstanding conjecture.

\clearpage
\setcounter{tocdepth}{1}
\hypersetup{bookmarksdepth=2}
\tableofcontents
\clearpage

\section{The repository question and the result obtained}
\label{sec:context}

\subsection{A precise boundary between existing work and this article}
The ProveIt transseries project studies formal asymptotic scales, well-based
supports, differential blocks, inversion, and analytic error transport.  At the
snapshot used here, the material has its own directory
\texttt{Analysis/Transseries/}.  The exact snapshot is
\begin{center}\small\ttfamily
a795fcffa4b3ece22761d81fbed3c43be8b81795.
\end{center}
The principal point of departure is the research package
\emph{A sharp regularity classification for countable exponential-feedback
transseries} \cite{feedback}.  Its README explicitly distinguishes the proved
left-sector Poincar\'e realization from a uniform Gevrey remainder theorem and
from directional Borel summation.  Those latter assertions were not established
there.  This article addresses that specific boundary.

The earlier package already proves the formal existence and Gevrey
classification of the feedback solution.  In particular, its classification
implies that the coefficient series is Gevrey--1 precisely at the growth
threshold $(j\log j)^2$.  We do not present that coefficient classification
as a discovery of the present paper.  Our main result is that, for this
nonnegative-slope model, the same threshold is exactly sufficient for
\emph{fine negative-direction Borel summation}, with identification of the
summed function.  We give an independent proof of the necessity needed here,
so our main theorem does not rely on accepting the earlier draft's asymptotic
proofs without review.

The focused source review covered the current project README, the
series-and-transseries directory README, and the feedback article and its
README.  It was not an exhaustive claim-by-claim audit of the much larger
canonical transseries volume.  The directory README also marks the recent
research arrivals as unmerged and not formalized \cite{project}.  The
formalization roadmap in Section~\ref{sec:lean} is accordingly prospective.

\subsection{Why inversion first changes the analytic problem}
Directly expanding $\U$ produces nonnegative, potentially very large
coefficients.  However, its inverse is characterized by
\begin{equation}
 u=\Phi(Q(u),u),\qquad
 \Phi(q,u)=\sum_{j\ge1}q^j\ee^{\lambda_j u}.
 \label{eq:literal}
\end{equation}
For $\Ree u\le0$, every coefficient $\ee^{\lambda_j u}$ has modulus at most
one.  The countably infinite kernel is therefore controlled by a geometric
series in $q$, regardless of how large the slopes are.  Inverting in $q$
while keeping $u$ fixed retains that information.  Expanding everything as
an ordinary power series at the outset discards it.

The resulting exact inverse has blocks of the form $u^k\ee^{bu}$.
Their integrated Borel transforms are Bessel kernels.  On the negative
Borel ray, positive $b$ produces oscillation instead of exponential growth;
the few possible negative frequencies satisfy a linear bound in $k$.
This is the mechanism behind the theorem.  It is a structural property of
the defining equation, not a general consequence of Gevrey regularity.

\subsection{Contribution map}
The new analytic assertions are the threshold equivalence of
Theorem~\ref{thm:threshold}, the explicit inverse Borel representation and
parabolic continuation of Theorem~\ref{thm:borel}, and the quadratic
remainder and error-certificate results of
Theorems~\ref{thm:remainder}--\ref{thm:optimal}.  The proof also gives an
all-slope continuous ray transform, even when no Borel germ exists.
Classical tools are proved in the specialized form needed here and credited
at their first use.  Numerical checks supplement, but do not establish,
the infinite-dimensional or asymptotic conclusions.

\section{Conventions and main theorems}
\label{sec:main}

Write
\begin{equation}
 \U(q)=\sum_{n\ge1}u_nq^n,
 \qquad \Q(u)=\sum_{n\ge1}v_nu^n,
 \qquad \Q=\U^{\inv}.
 \label{eq:formal}
\end{equation}
The formal inverse exists because $u_1=1$.  All slopes are finite real
numbers, and are nonnegative throughout.  Natural logarithms are used.
A statement about $j\log j$ starts at $j=2$.

\begin{definition}[Gevrey--1 and integrated Borel transform]
A formal series $\widehat f(u)=\sum_{n\ge0}f_nu^n$ is in $\G$ when
$|f_n|\le C H^n n!$ for some $C,H>0$.  Its integrated Borel transform is
\begin{equation}
 \B\widehat f(\zeta)=\sum_{n\ge0}\frac{f_n}{n!}\zeta^n.
 \label{eq:B0}
\end{equation}
Thus $\widehat f\in\G$ exactly when \eqref{eq:B0} is a holomorphic germ at
zero.  The usual Borel minor of a series with zero constant term is the
derivative of \eqref{eq:B0}.
\end{definition}

\begin{definition}[Fine summability on the negative ray]
\label{def:fine}
Let
\[
 \Omega^-_\delta=\{\zeta\in\C:\dist(\zeta,(-\infty,0])<\delta\}.
\]
A series is \emph{finely Borel summable in direction $\pi$} when its Borel
germ extends to some $\Omega^-_\delta$ and, there, satisfies
$|\B\widehat f(\zeta)|\le C\ee^{H|\zeta|}$.
Its sum at $u=-r$, for sufficiently small $r>0$, is
\begin{equation}
 \mathcal S_\pi\widehat f(-r)
  =\frac1r\int_0^\infty\ee^{-s/r}\B\widehat f(-s)\,\dd s.
 \label{eq:laplace}
\end{equation}
For complex $u$ the same formula is
$-u^{-1}\int_0^\infty \ee^{s/u}\B\widehat f(-s)\,\dd s$, on
$\Ree(-1/u)>H$.
\end{definition}

This is the fine, half-strip version of summability, as distinguished from
summability in an open arc of directions in \cite[Sections 7 and 9]{sauzin}.
Using the Borel minor gives an equivalent definition: integration preserves
exponential growth, and differentiation preserves it after narrowing the
strip.  The domain in the $u$-plane is a disc tangent to the origin from
the left, not a sector of opening greater than $\pi$.

\begin{theorem}[Sharp negative-direction threshold]
\label{thm:threshold}
For every sequence $\lambda_j\ge0$, the following five conditions are
equivalent:
\begin{enumerate}[label=(\roman*)]
\item $\lambda_j=O((j\log(j+1))^2)$;
\item $\U\in\G$;
\item $\Q\in\G$;
\item $\U$ is finely Borel summable in direction $\pi$;
\item $\Q$ is finely Borel summable in direction $\pi$.
\end{enumerate}
When they hold, the summed functions are compositional inverses near the
negative real axis.  In particular, for small negative $q$,
\begin{equation}
 \mathcal S_\pi\U(q)
 =\sum_{j\ge1}q^j\exp\!\left(\lambda_j\mathcal S_\pi\U(q)\right),
 \label{eq:summed-feedback}
\end{equation}
and the sum on the right converges absolutely.  The analogous identity
holds on a sufficiently small common complex neighborhood contained in
the left sectors where the dependent variable has negative real part.
\end{theorem}

\begin{theorem}[Constructive inverse continuation]
\label{thm:borel}
Set
\begin{equation}
 A=\limsup_{j\to\infty}\frac{\lambda_j}{(j\log j)^2}.
 \label{eq:A}
\end{equation}
If $0<A<\infty$, the explicit kernel series of
Section~\ref{sec:bessel} defines $\B\Q$ holomorphically on
\begin{equation}
 \mathcal P_A=\left\{\zeta\in\C:
       |\zeta|+\Ree\zeta<\frac1{2A}\right\}.
 \label{eq:parabola}
\end{equation}
It has exponential growth on a sufficiently thin fixed tubular
neighborhood of the negative ray.  If $A=0$, the same kernel series is
entire and still has this negative-ray growth property.  The set
\eqref{eq:parabola} is a guaranteed continuation region, not a claim that
its boundary consists of singularities.
\end{theorem}

The quadratic case has a stronger uniform conclusion.  It is convenient
to retain the convergent factor $\ee^{-u}$ exactly and write
\begin{equation}
 P(u)=\ee^uQ(u),\qquad \Pp(u)=\sum_{n\ge1}a_nu^n.
 \label{eq:P}
\end{equation}
Let $\rho=3-2\sqrt2$ and let $\sh_k$ be defined in
\eqref{eq:schroeder} below.

\begin{theorem}[Uniform logarithmic-Gevrey remainder]
\label{thm:remainder}
Suppose $\lambda_j=j^2$.  Boundary values refer to the continuous
block-series extension constructed below.  For $0<r_0<\rho$, $N\ge2$, and
$\Ree u\le0$, $|u|=r\le r_0$, one has
\begin{equation}
 \left|P(u)-\sum_{n=1}^{N-1}a_nu^n\right|
 \le r^N M_N(r_0),
 \label{eq:remainder}
\end{equation}
where the entirely explicit positive majorant is
\begin{equation}
 M_N(r_0)=\frac{\rho^{-N}}{1-r_0/\rho}
 +\sum_{k=1}^{N-1}\sh_k
       \frac{[k(k-1)]^{N-k}}{(N-k)!}.
 \label{eq:MN}
\end{equation}
For fixed $r_0$,
\begin{equation}
 \log M_N(r_0)
 =N\log N-2N\log\log N+(\log4-1)N+o(N).
 \label{eq:MN-asymp}
\end{equation}
The error for approximating $Q(u)$ by
$\ee^{-u}\sum_{n<N}a_nu^n$ is at most $\ee^r$ times the right side of
\eqref{eq:remainder}.
\end{theorem}

\begin{theorem}[A certified beyond-all-orders scale]
\label{thm:optimal}
In the quadratic case, put $L=\log(1/r)$ and
$N(r)=\lfloor L^2/(4r)\rfloor$.  For every $\varepsilon>0$ and all
sufficiently small $r>0$,
\begin{equation}
 \sup_{\substack{|u|=r\\\Ree u\le0}}
 \left|P(u)-\sum_{n<N(r)}a_nu^n\right|
 \le \exp\!\left[-\left(\frac14-\varepsilon\right)
                         \frac{L^2}{r}\right].
 \label{eq:optimal}
\end{equation}
The same conclusion holds for the exponentially normalized approximation
to $Q$, after possibly decreasing the admissible $r$.
\end{theorem}

Here and below, ``optimal'' in a theorem label refers to the leading-order
optimization of our proved majorant.  It does not assert a matching lower
bound for the actual truncation error.  Cancellation can make that error
smaller.  No Stokes constant is inferred from the number $1/4$.
% ed. (2026-09-29): scope note on thm:optimal (an upper bound from the majorant, not a sharp optimal-truncation result).
\begin{ednote}
As the preceding paragraph states, Theorem~\ref{thm:optimal} is an upper
bound derived from the majorant $M_N$; despite its label it is not a sharp
optimal-truncation theorem, and the repository does not count it among
sharp instances of optimal truncation. Later packages in
\nolinkurl{Analysis/Transseries/docs/series-and-transseries/} meet the same
scale for other quantities:
\nolinkurl{Microscopic_Condensation_Exponential_Feedback/} (its
Theorem~\texttt{thm:least}) gives the value of the least term of the \emph{forward}
series, $\exp(-(1+o(1))L^2/(4t))$ at an index $\sim L^2/(4t)$ with
$L=\log(1/t)$, and
\nolinkurl{Sharp_Negative_Direction_Summability_Exponential_Feedback/}
(its equation \texttt{eq:quadraticactual}) bounds the actual truncation error of
$U$ on proper left sectors by $C\exp[-c\log^2(1/|q|)/|q|]$ with an
unspecified constant. No matching lower bound for the actual error is
known.
\end{ednote}

\section{Formal existence and the necessary growth condition}
\label{sec:formal}

\subsection{Local finiteness and positive Lagrange weights}
Although $\Phi$ need not be analytic near $(0,0)$, it is legitimate
formally: its coefficient of $q^n$ uses only $j\le n$.  If
$V,W\in q\R[[q]]$ agree through degree $n-1$, then
$\Phi(q,V)$ and $\Phi(q,W)$ agree through degree $n$.  Iteration starting
at zero therefore converges coefficientwise to a unique fixed point.
Its first coefficient is one, and every coefficient is nonnegative.

For completeness, classical Lagrange inversion \cite{gessel} yields a
particularly useful finite coefficient formula.  Introduce an auxiliary
parameter $t$ in $V=t\Phi(q,V)$, invert with respect to $t$, and then set
$t=1$ coefficientwise.  Each coefficient of $q$ permits only finitely many
powers of $t$.  For a finitely supported sequence of nonnegative integers
$\mathbf m=(m_1,m_2,\ldots)$, set
$|\mathbf m|=\sum m_j$ and $w(\mathbf m)=\sum jm_j$.  The result is
\begin{equation}
 u_n=\sum_{w(\mathbf m)=n}
       \frac{(\sum_j\lambda_jm_j)^{|\mathbf m|-1}}
            {\prod_jm_j!}.
 \label{eq:positive-Lagrange}
\end{equation}
The convention $0^0=1$ handles single-component terms.  This identity is
only a finite formal calculation; it imposes no analytic condition on
the sequence of slopes.

Taking one component of size $j\ge2$ and $m$ components of size one in
\eqref{eq:positive-Lagrange} gives
\begin{equation}
 u_{j+m}\ge \frac{(\lambda_j+m\lambda_1)^m}{m!}
             \ge\frac{\lambda_j^m}{m!}.
 \label{eq:lower}
\end{equation}
This elementary inequality already detects the correct necessary scale.

\begin{lemma}[Necessity of the logarithmic-square threshold]
\label{lem:necessary}
If $\U\in\G$, then $\lambda_j=O((j\log(j+1))^2)$.
\end{lemma}
\begin{proof}
Choose $C,H\ge1$ with $u_n\le CH^n n!$, and put
$m=\lceil j\log(j+1)\rceil$.  From \eqref{eq:lower},
\[
 \lambda_j\le C^{1/m}H^{1+j/m}
                   \bigl(m!(j+m)!\bigr)^{1/m}.
\]
Using $n!\le n^n$, the last factor is at most
$m(j+m)^{1+j/m}$.  Now $j/m=O(1/\log j)$, so
$(j+m)^{j/m}$ stays bounded, while $m$ and $j+m$ are both
$O(j\log(j+1))$.  The remaining factors are bounded as well.  The
finitely many small $j$ are absorbed into the constant.
\end{proof}

\subsection{Why reversion does not change class membership}
\begin{lemma}[Gevrey--1 reversion]
\label{lem:gevreyinverse}
Let $f(z)=z+\sum_{n\ge2}f_nz^n$ and let $g=f^{\inv}$ be formal.
Then $f\in\G$ if and only if $g\in\G$.
\end{lemma}
\begin{proof}
It suffices to prove one direction, since $f=g^{\inv}$.  Lagrange's
coefficient formula is
\[
 g_n=\frac1n\coeff{z}{n-1}
       \left(1+\sum_{j\ge1}f_{j+1}z^j\right)^{-n}.
\]
Expand the negative power.  A term with $\ell$ factors is indexed by a
composition $j_1+\cdots+j_\ell=n-1$, with $j_i\ge1$, and has a binomial
factor $\binom{n+\ell-1}{\ell}$.  If
$|f_m|\le CH^m m!$, the product of factorials in that term obeys
\begin{equation}
 \prod_{i=1}^{\ell}(j_i+1)!
 \le 2^{\ell-1}(n-\ell+1)!\le 2^n n!.
 \label{eq:factorial-product}
\end{equation}
To see the first inequality, move a unit from the smaller of two parts
greater than one to the larger; the factorial product can only increase.
Iterating leaves one large part and $\ell-1$ parts equal to one.

The binomial factor is at most $4^n$, the number of all relevant
compositions is at most $2^n$, and the factors involving $C$ and $H$ are
at most $D^n$ for a fixed $D$.  Summing gives $|g_n|\le E^n n!$ after
enlarging $E$.  This proves class membership, not preservation of the
sharp factorial-normalized type.
\end{proof}

The preceding lemma is included to make the logic of necessity independent
of any sharp asymptotic classification.  Fine summability implies Gevrey--1
by local analyticity of the Borel germ.  Thus either fine summability
assertion in Theorem~\ref{thm:threshold} already forces the slope condition.
The rest of the paper proves its sufficiency constructively.

\section{An exact signed exponential-block inverse}
\label{sec:blocks}

\subsection{The coefficient measures}
Fix $u$ temporarily and invert the map $q\mapsto\Phi(q,u)$ at zero, with
an independent target variable $w$.  Lagrange inversion gives
\[
 q(w;u)=\sum_{k\ge1}\frac{w^k}{k}
       \coeff{q}{k-1}
       \left(\sum_{j\ge1}q^{j-1}\ee^{\lambda_j u}\right)^{-k}.
\]
Now set $w=u$.  For $k\ge2$, let $1\le\ell\le k-1$ and let
$\mathbf m=(m_1,\ldots,m_\ell)$ range over ordered positive compositions
of $k-1$.  Define
\begin{align}
 \omega_{k,\ell}&=\frac{(-1)^\ell}{k}
                         \binom{k+\ell-1}{\ell},\label{eq:weight}\\
 b_{k,\mathbf m}&=\sum_{i=1}^\ell\lambda_{m_i+1}
                              -(k+\ell)\lambda_1.\label{eq:frequency}
\end{align}
Let $\mu_k$ be the finite signed atomic measure placing mass
$\omega_{k,\ell}$ at $b_{k,\mathbf m}$ for each indexed composition.
Coincident atoms are added.  Put $\mu_1=\delta_{-\lambda_1}$.  Formally,
\begin{equation}
 \boxed{\quad Q(u)=\sum_{k\ge1}u^k\int_\R\ee^{bu}\,\dd\mu_k(b).\quad}
 \label{eq:block-expansion}
\end{equation}
Every ordinary Taylor coefficient uses only $k\le n$, so
\eqref{eq:block-expansion} also defines a formal series without convergence
hypotheses.  It solves $u=\Phi(Q(u),u)$ formally.  Uniqueness of the
formal feedback fixed point after substitution $q=Q(u)$ gives
$\U(Q(u))=u$, identifying this series with the compositional inverse
in \eqref{eq:formal}.

\subsection{A universal combinatorial majorant}
Define $\sh_1=1$ and, for $k\ge2$,
\begin{equation}
 \sh_k=\frac1k\sum_{\ell=1}^{k-1}
                \binom{k+\ell-1}{\ell}\binom{k-2}{\ell-1}.
 \label{eq:schroeder}
\end{equation}
There are $\binom{k-2}{\ell-1}$ relevant compositions, so
$\|\mu_k\|_{\TV}\le\sh_k$.

The first seven values are $1,1,3,11,45,197,903$.

\begin{lemma}[Generating function and usable bounds]
\label{lem:mass}
The generating function of these positive majorants is
\begin{equation}
 S(t)=\sum_{k\ge1}\sh_kt^k
     =\frac{1+t-\sqrt{1-6t+t^2}}4.
 \label{eq:S}
\end{equation}
Its radius of convergence is $\rho=3-2\sqrt2$, and
\begin{equation}
 \sh_k\le\rho^{-k},\qquad \sh_k\le\frac{6^k}{4}.
 \label{eq:C-bounds}
\end{equation}
\end{lemma}
\begin{proof}
Applying the coefficient form of Lagrange inversion to
$S=t(1-S)/(1-2S)$ gives exactly \eqref{eq:schroeder}.  The branch with
$S(0)=0$ of $2S^2-(1+t)S+t=0$ is \eqref{eq:S}.  Its first positive
branch point is $\rho$.  Since all coefficients are nonnegative,
$\sh_k\rho^k\le S(\rho)=(1+\rho)/4<1$; convergence at $\rho$ follows
also by monotone passage in \eqref{eq:S}.  Similarly
$\sh_k6^{-k}\le S(1/6)=1/4$.
\end{proof}

These are the small Schr\"oder coefficients, but no enumerative
interpretation is required for any subsequent argument.  Retaining the
uncancelled mass majorant, rather than the actual total variation after
coincident atoms cancel, makes every estimate transparent and reproducible.

\begin{proposition}[All-slope analytic inverse]
\label{prop:all-slopes}
For every nonnegative slope sequence, \eqref{eq:block-expansion} converges
normally in a sufficiently small open left half-disc and uniformly on
smaller closed left half-discs.  More explicitly, it is absolutely
convergent whenever
\begin{equation}
 \Ree u\le0,\qquad |u|\ee^{2\lambda_1|u|}<\rho.
 \label{eq:block-domain}
\end{equation}
The resulting function satisfies \eqref{eq:literal} for sufficiently small
$u$ in that half-disc.  It has $\Q$ as a uniform Poincar\'e expansion on
each smaller closed left half-disc, as $u\to0$.
\end{proposition}
\begin{proof}
Nonnegativity of the slopes gives $b\ge-2k\lambda_1$ at every atom.
For $\Ree u\le0$,
\[
 \left|u^k\int\ee^{bu}\,\dd\mu_k(b)\right|
 \le\sh_k\bigl(|u|\ee^{2\lambda_1|u|}\bigr)^k.
\]
This proves normal and boundary-uniform convergence.

For a fixed $u$ with negative real part,
$q\mapsto\Phi(q,u)$ is holomorphic for $|q|<1$ and has nonzero derivative
$\ee^{\lambda_1u}$ at zero.  Its local inverse in $w$ has the Lagrange
coefficients used above.  The same majorant shows convergence of that
inverse for small $w$, uniformly for small $u$ in the left half-disc.
For $w=u$ both the inverse series and its value remain sufficiently small;
substitution in the convergent $q$-series gives \eqref{eq:literal}.
This justifies the functional identity, not only the formal coefficients.

For each fixed $N$, the blocks with $k<N$ are finite sums of entire
exponentials and can be Taylor expanded with a remainder of order $u^N$.
The blocks with $k\ge N$ have a geometric $O(|u|^N)$ tail on smaller
half-discs.  The constants may depend arbitrarily on $N$ and the slopes.
This proves uniform Poincar\'e asymptotics, without a Gevrey assertion.
\end{proof}

\section{Bessel kernels and a ray transform for arbitrary slopes}
\label{sec:bessel}

\subsection{One block under the Borel transform}
For $k\ge1$ and $b\in\R$ define the entire function
\begin{equation}
 K_{k,b}(\zeta)
  =\sum_{m\ge0}\frac{b^m\zeta^{m+k}}{m!(m+k)!}
  =\frac{\zeta^k}{k!}\,{}_0F_1(k+1;b\zeta).
 \label{eq:kernel}
\end{equation}
It is the integrated Borel transform of $u^k\ee^{bu}$.
A normalized form of the classical Poisson--Bessel integral
\cite[Eq.~10.9.4]{dlmf} is particularly effective here.

\begin{lemma}[Probability-integral representation]
\label{lem:kernel}
Let
\[
 p_k(v)=\frac{\Gamma(k+1)}{\sqrt\pi\Gamma(k+1/2)}
                    (1-v^2)^{k-1/2},\qquad -1<v<1.
\]
Then $\int_{-1}^1p_k(v)\,\dd v=1$ and
\begin{equation}
 K_{k,b}(\zeta)=\frac{\zeta^k}{k!}
        \int_{-1}^1\exp(2\sqrt{b\zeta}\,v)p_k(v)\,\dd v.
 \label{eq:poisson}
\end{equation}
The right side is independent of the choice of square root.  If $b\ge0$,
\begin{equation}
 |K_{k,b}(\zeta)|\le\frac{|\zeta|^k}{k!}
                   \exp(2\sqrt b\,|\Ree\sqrt\zeta|).
 \label{eq:pos-kernel}
\end{equation}
If $b<0$, replace $\sqrt b\,|\Ree\sqrt\zeta|$ by
$\sqrt{|b|}\,|\Imm\sqrt\zeta|$ in that bound.  In particular, for $s\ge0$,
\begin{align}
 |K_{k,b}(-s)|&\le s^k/k!,&& b\ge0,\label{eq:osc-bound}\\
 |K_{k,b}(-s)|&\le(s^k/k!)\ee^{2\sqrt{|b|s}},&&b<0.\label{eq:neg-bound}
\end{align}
For $b>0$ there is also the explicit oscillatory formula
\begin{equation}
 K_{k,b}(-s)=(-1)^k(s/b)^{k/2}J_k(2\sqrt{bs}).
 \label{eq:besselJ}
\end{equation}
\end{lemma}
\begin{proof}
A beta integral gives the normalization and the even moments
\[
 \int_{-1}^1v^{2m}p_k(v)\,\dd v
  =\frac{(2m)!k!}{4^m m!(m+k)!}.
\]
Odd moments vanish.  Expanding the exponential in \eqref{eq:poisson}
and using these moments gives \eqref{eq:kernel}.  The same evenness proves
square-root independence.  Taking absolute values under a probability
integral proves \eqref{eq:pos-kernel} and its negative-$b$ counterpart.
For a negative argument and positive $b$, the integrand is oscillatory
and has modulus one.  Formula \eqref{eq:besselJ} follows by comparing the
power series for $J_k$.
\end{proof}

\subsection{Exact Laplace recovery and the distinction from a Borel germ}
\begin{lemma}[One-block recovery]
\label{lem:kernel-laplace}
For every $r>0$, $k\ge1$, and real $b$,
\begin{equation}
 \frac1r\int_0^\infty\ee^{-s/r}K_{k,b}(-s)\,\dd s
       =(-r)^k\ee^{-br}.
 \label{eq:kernel-laplace}
\end{equation}
Moreover, the same integral with an absolute value inside is at most
$r^k$ for $b\ge0$, and equals $r^k\ee^{|b|r}$ for $b<0$.
\end{lemma}
\begin{proof}
For one fixed $b$, termwise integration of \eqref{eq:kernel} is justified
absolutely: the integrated series of absolute values is
$r^k\sum_{m\ge0}(|b|r)^m/m!$.  The signed sum is the right side of
\eqref{eq:kernel-laplace}.  For $b\ge0$, use the stronger oscillatory
bound \eqref{eq:osc-bound} before integrating.  For $b<0$, every term of
$(-1)^kK_{k,b}(-s)$ is nonnegative, so the absolute integral is the
positive exponential sum just computed.
\end{proof}

\begin{proposition}[Continuous ray representation without a growth hypothesis]
\label{prop:ray}
For arbitrary nonnegative slopes, the function
\begin{equation}
 \mathfrak B_Q(-s)=\sum_{k\ge1}\int K_{k,b}(-s)\,\dd\mu_k(b),
 \qquad s\ge0,
 \label{eq:ray-series}
\end{equation}
is well-defined, continuous, and at most exponentially growing in $s$.
For all sufficiently small $r>0$ its Laplace integral recovers the
analytic inverse of Proposition~\ref{prop:all-slopes}:
\begin{equation}
 Q(-r)=\frac1r\int_0^\infty\ee^{-s/r}\mathfrak B_Q(-s)\,\dd s.
 \label{eq:ray-recovery}
\end{equation}
This alone is \emph{not} a fine Borel-summability assertion: without a
slope bound, \eqref{eq:ray-series} need not be the restriction of a
holomorphic Borel germ at zero.
\end{proposition}
\begin{proof}
Since a negative frequency has $|b|\le2k\lambda_1$, the $k$th term is
bounded on the ray by
\[
 \sh_k\frac{s^k}{k!}\exp(2\sqrt{2\lambda_1ks}).
\]
For any fixed $\eta>0$, Young's inequality bounds the last exponential
by $\exp(\eta k+2\lambda_1s/\eta)$.  Sum with
$\sh_k\le\rho^{-k}$.  The result is an exponential in $s$, and the same
bound gives uniform convergence on compact ray intervals.

Lemma~\ref{lem:kernel-laplace} bounds the integrated absolute contribution
of block $k$ by
$\sh_kr^k\ee^{2k\lambda_1r}$.  Its sum is finite under
\eqref{eq:block-domain}.  Fubini's theorem is therefore applicable to
the kernel sum and Laplace integral; it gives \eqref{eq:block-expansion}
at $u=-r$ exactly.  No exchange of the original, generally divergent,
forward power series with a Laplace integral has been made.
\end{proof}

The last warning is substantive.  A continuous ray function with a
formal expansion can be produced even when that formal expansion has
zero Borel radius.  Fine summability requires the additional complex
analytic neighborhood proved next.

\section{The parabolic continuation and fine growth estimate}
\label{sec:continuation}

\subsection{From a slope bound to a block-frequency bound}
\begin{lemma}[Concentration of the maximal frequency]
\label{lem:frequency}
If $A<\infty$ in \eqref{eq:A}, then for every $a>A$ there is $D\ge0$
such that every positive frequency in block $k\ge2$ satisfies
\begin{equation}
 b\le a(k\log k)^2+D(k-1).
 \label{eq:b-bound}
\end{equation}
Every negative frequency satisfies $|b|\le2k\lambda_1$.
\end{lemma}
\begin{proof}
Choose $D$ so that, for all $j\ge2$,
$\lambda_j\le a(j\log j)^2+D(j-1)$.  Such a $D$ exists by the definition
of the limsup and the finiteness of each exceptional slope.  The function
\[
 g(x)=((x+1)\log(x+1))^2,\qquad x\ge0,
\]
is convex and has $g(0)=0$.  Indeed, for $t\ge1$ the second derivative of
$t^2\log^2t$ is $2\log^2t+6\log t+2>0$.  Convexity and $g(0)=0$ imply
$g(x)+g(y)\le g(x+y)$ for nonnegative $x,y$: bound each term by its
fraction of the chord from zero to $x+y$.  Therefore
\[
 \sum_{i=1}^\ell\lambda_{m_i+1}
 \le a\,g\!\left(\sum_i m_i\right)+D\sum_i m_i
 =a(k\log k)^2+D(k-1).
\]
Dropping the nonpositive subtraction in \eqref{eq:frequency} proves
\eqref{eq:b-bound}.  The negative-frequency assertion was already used
in Proposition~\ref{prop:all-slopes}.
\end{proof}

\subsection{Normal convergence in the Borel plane}
Define
\begin{equation}
 B_Q(\zeta)=\sum_{k\ge1}\int K_{k,b}(\zeta)\,\dd\mu_k(b).
 \label{eq:BQ}
\end{equation}
The square-root quantities
\[
 h(\zeta)=|\Ree\sqrt\zeta|,
 \qquad \widetilde h(\zeta)=|\Imm\sqrt\zeta|
\]
are independent of the sign of the root, and
$h(\zeta)^2=(|\zeta|+\Ree\zeta)/2$.
By Lemmas~\ref{lem:kernel} and \ref{lem:frequency}, the absolute mass of
block $k\ge2$ is bounded by
\begin{equation}
 \frac{\rho^{-k}|\zeta|^k}{k!}
 \exp\!\left(2\sqrt a\,h(\zeta)k\log k
       +2\sqrt D\,h(\zeta)\sqrt k
       +2\sqrt{2\lambda_1}\,\widetilde h(\zeta)\sqrt k\right).
 \label{eq:master-bound}
\end{equation}
The estimate combines the positive- and negative-frequency bounds; an
individual atom needs only one of the corresponding exponents.

Let $K$ be compact and suppose $\sup_K2\sqrt a\,h<1$.  Taking logarithms
in \eqref{eq:master-bound} and using Stirling's elementary estimate shows
that the $k$th bound is
\[
 \exp\bigl(-(1-\beta)k\log k+O_K(k)+O_K(\sqrt k)\bigr),
 \qquad \beta=\sup_K2\sqrt a\,h<1.
\]
This is summable.  Thus \eqref{eq:BQ} converges normally on
$\{2\sqrt a\,h<1\}$.  If $A>0$, every compact subset of
\eqref{eq:parabola} admits a choice $a>A$ satisfying this strict
inequality.  If $A=0$, every compact subset of $\C$ admits a sufficiently
small $a>0$.  We obtain the continuation assertions of
Theorem~\ref{thm:borel}.

In a neighborhood of zero, normal convergence also permits Taylor
coefficient extraction.  Each kernel in block $k$ vanishes to order $k$;
hence the coefficient of $\zeta^n$ is a finite sum and is exactly $v_n/n!$.
Consequently \eqref{eq:BQ} really is the analytic continuation of
$\B\Q$, not merely a ray function chosen to have a desired Laplace
transform.

\subsection{Exponential bounds in a fixed strip}
Parabolic continuation alone does not prove fine summability; growth
must be controlled on a fixed-width neighborhood of the entire ray.
We give that estimate explicitly.

Choose one $a>A$ and choose $\delta>0$ small enough that the closure of
$\Omega^-_{2\delta}$ stays inside $\{2\sqrt a\,h<1\}$.  This is possible
because $h^2\le2\delta$ on that tube.  For
$\zeta=-s+it$, $s\ge1$, $|t|\le\delta$, one has
\begin{equation}
 h(\zeta)\le\frac{\delta}{2\sqrt s},\qquad
 \widetilde h(\zeta)\le\sqrt{s+\delta},\qquad
 |\zeta|\le s+\delta.
 \label{eq:strip-roots}
\end{equation}
The first estimate follows by rationalizing
$\sqrt{s^2+t^2}-s$.

Absorbing fixed constants into $d_1,d_2,d_3$, the logarithm of the
$k$th bound in \eqref{eq:master-bound} is at most
\begin{equation}
 k\log\bigl(\rho^{-1}(s+\delta)\bigr)-\log k!
 +\frac{d_1}{\sqrt s}k\log k
 +\frac{d_2}{\sqrt s}\sqrt k+d_3\sqrt{ks}.
 \label{eq:strip-term}
\end{equation}
Split the sum into $k\le s^2$ and $k>s^2$.

In the first range,
$(d_1/\sqrt s)\log k\le2d_1\log s/\sqrt s=o(1)$.
Young's inequality bounds the last two terms by $c_1k+c_2s+c_3$.
For large $s$, the first part of the sum is consequently at most
\[
 C\ee^{c_2s}\sum_{k\ge1}
       \frac{(\rho^{-1}\ee^{c_1+1}(s+\delta))^k}{k!}
 \le C'\ee^{H_1s}.
\]
In the second range, $\log(s+\delta)\le\tfrac12\log k+o(1)$.
For sufficiently large $s$ we have $d_1/\sqrt s\le1/8$.
Using $\log k!\ge k\log k-k$ in \eqref{eq:strip-term}, while
$\sqrt{ks}/k\le1/\sqrt s$, gives an upper bound
$-\tfrac14 k\log k$ after increasing the threshold on $s$.
The tail is thus uniformly bounded.  The remaining bounded portion of
the tube is compactly contained in the domain of normal convergence.
We conclude that
\begin{equation}
 |B_Q(\zeta)|\le C\ee^{H|\zeta|},\qquad \zeta\in\Omega^-_\delta.
 \label{eq:fine-growth}
\end{equation}
This proves the remaining assertion of Theorem~\ref{thm:borel} and, in
particular, the implication (i)$\Rightarrow$(v) of
Theorem~\ref{thm:threshold}.  Proposition~\ref{prop:ray} then identifies
the Borel sum with the literal analytic inverse $Q$.

\section{Summation commutes with taking the inverse}
\label{sec:nonlinear}

The preceding construction proves fine summability of $\Q$ directly.
To transfer the result to $\U$, one needs closure under compositional
inversion, not just closure of Gevrey coefficient classes.  This is a
classical nonlinear summation principle.  We include a proof adapted to
fixed half-strips, following the convolution method underlying
\cite[Sections 13 and 17]{sauzin}.  This avoids silently replacing a
half-strip hypothesis by an open angular-sector hypothesis.

\subsection{The elementary half-strip algebra}
At infinity write
$\varphi(z)=\sum_{n\ge1}\varphi_nz^{-n}$ and use its Borel minor
\[
 \mathscr B\varphi(\xi)
    =\sum_{n\ge1}\frac{\varphi_n}{(n-1)!}\xi^{n-1}.
\]
Let $\Omega^+_\delta=-\Omega^-_\delta$.  Consider those minors holomorphic
there and bounded by $M\ee^{c|\xi|}$.  The tube is convex, so every straight
segment from zero to a point of the tube stays inside it.  If $\widehat
\varphi$ is such a minor, its $k$-fold convolution obeys
\begin{equation}
 |\widehat\varphi^{*k}(\xi)|
 \le M^k\frac{|\xi|^{k-1}}{(k-1)!}\ee^{c|\xi|}.
 \label{eq:conv-bound}
\end{equation}
Indeed, parameterize the convolution simplex by nonnegative coordinates
summing to one; the exponential factors multiply to
$\ee^{c|\xi|}$ and the simplex volume is $1/(k-1)!$.

Products of formal series correspond to convolution, with the constant
terms handled separately.  Formula \eqref{eq:conv-bound} also shows that
substitution into a convergent scalar power series is permitted whenever
the substituted series has zero constant term.  If $H(w)=\sum h_kw^k$
and $|h_k|\le CR^{-k}$, the sum of the Borel minors of
$\sum_{k\ge1}h_k\varphi^k$ is bounded by a constant times
$\exp((c+M/R)|\xi|)$.  In particular, reciprocals of series with nonzero
constant term belong to the same class.

Multiplication by $z$, with the new polynomial part separated off,
corresponds to differentiation of the minor.  Cauchy's estimate on a
narrower tube preserves an exponential bound.  Division by $z$
corresponds to integration.  Translation $z\mapsto z+c_0$ multiplies
the minor by $\ee^{-c_0\xi}$ and preserves the class.  Laplace summation
commutes with all these operations: products follow from Fubini on
convolution integrals, reciprocal and scalar substitution from the same
absolutely convergent convolution majorant, multiplication by $z$ from
integration by parts, and translation from its exponential multiplier.
All these identities hold on a sufficiently far right half-plane.

\subsection{A direct inverse formula and its analytic meaning}
\begin{lemma}[Fine-summable inverse at infinity]
\label{lem:fineinverse}
Suppose $f(z)=z+c_0+\varphi(z)$, where $\varphi\in z^{-1}\C[[z^{-1}]]$
has a fine-summable minor in the positive direction.  Then the formal
inverse $g(z)=z-c_0+O(z^{-1})$ has the same property.  Its sum is the
analytic inverse of the sum of $f$ on sufficiently far right half-planes.
\end{lemma}
\begin{proof}
First take $c_0=0$.  Formal Lagrange inversion gives
\begin{equation}
 g(z)=z+\chi(z),\qquad
 \chi=\sum_{k\ge1}\frac{(-1)^k}{k!}
                       \partial_z^{\,k-1}(\varphi^k).
 \label{eq:inverse-operator}
\end{equation}
One way to verify this formula is to solve
$\xi=-t\varphi(z+\xi)$ in the formal parameter $t$ and then put $t=1$;
the order in $z^{-1}$ tends to infinity with $k$.
Since differentiation corresponds to multiplication by $-\xi$ in the
Borel plane, the candidate minor is
\begin{equation}
 \widehat\chi(\xi)
 =-\sum_{k\ge1}\frac{\xi^{k-1}}{k!}
                          \widehat\varphi^{*k}(\xi).
 \label{eq:inverse-minor}
\end{equation}
By \eqref{eq:conv-bound}, its absolute value is at most
\[
 \ee^{c|\xi|}\sum_{k\ge1}
        \frac{M^k|\xi|^{2k-2}}{k!(k-1)!}
 \le M\exp\bigl((c+2\sqrt M)|\xi|\bigr).
\]
The last inequality follows by comparison with
$\sum_{j\ge0}(\sqrt M|\xi|)^{2j}/(j!)^2\le
\exp(2\sqrt M|\xi|)$.  Normal convergence proves the required
holomorphy on the same half-strip.

It remains to identify the sum as an inverse function.  The Laplace
sum $\varphi_*(z)$ satisfies
$|\varphi_*(z)|\le M/(\Ree z-c)$ and
$|\varphi_*'(z)|\le M/(\Ree z-c)^2$.
For $\Ree z$ sufficiently large and $|t|\le2$, the equation
$\xi=-t\varphi_*(z+\xi)$ is a contraction on $|\xi|\le1$.
Its solution is holomorphic in $t$ on a disc of radius greater than one.
The ordinary analytic Lagrange formula is therefore convergent at $t=1$
and gives \eqref{eq:inverse-operator} with $\varphi_*$ in place of
$\varphi$.  On the other hand, the exponential majorant for
\eqref{eq:inverse-minor} permits termwise Laplace integration, and each
term sums to the corresponding term of that analytic Lagrange formula.
Thus the sum of $g$ really is the inverse of $z+\varphi_*(z)$.

For a nonzero constant, write $f=\tau_{c_0}\circ f_0$, where
$f_0=z+\varphi$ and $\tau_{c_0}(z)=z+c_0$.
Then $f^{\inv}=f_0^{\inv}\circ\tau_{-c_0}$.  Translation, already
proved to preserve fine summability and its sums, completes the argument.
\end{proof}

\subsection{Return to the feedback variables}
For a tangent-to-identity series $Q$, conjugate by the involution
$J(z)=-1/z$:
\begin{equation}
 f=J\circ Q\circ J,
 \qquad f(z)=-\frac1{Q(-1/z)}=z+c_0+O(z^{-1}).
 \label{eq:conjugation}
\end{equation}
Fine summability of $Q$ on the negative ray gives the positive-direction
half-strip minor of $Q(-1/z)$ by reflection and differentiation of
$B_Q(-\xi)$.  The algebra operations just proved, including multiplication
by $z$ and reciprocal of a nonzero constant term, give the fine-summable
tail of $f$.  Lemma~\ref{lem:fineinverse} gives $g=f^{\inv}$, and
conjugating back gives
$J\circ g\circ J=Q^{\inv}=U$.  Each operation respects the summed
functions.  We have proved fine summability of $U$, and that the two sums
are inverses, without assuming a larger angular continuation domain.

\begin{proof}[Completion of Theorem~\ref{thm:threshold}]
The slope bound implies fine summability of $Q$ by
Theorem~\ref{thm:borel}, and the preceding argument implies fine
summability of $U$.  Fine summability implies Gevrey--1.  The two Gevrey
conditions are equivalent by Lemma~\ref{lem:gevreyinverse}, and either
implies the slope condition by Lemma~\ref{lem:necessary}.
This proves all five equivalences.

The sum of $Q$ is the literal inverse constructed in
Proposition~\ref{prop:all-slopes}, by exact Laplace recovery.  For small
negative $q$, its inverse satisfies $U(q)=q+O(q^2)<0$.
Putting $u=U(q)$ in \eqref{eq:literal} gives
\eqref{eq:summed-feedback}; its absolute convergence follows from
$\Ree u<0$ and $|q|<1$.  On a sufficiently small proper left sector,
$U(q)=q+O(q^2)$ similarly implies $\Ree U(q)<0$, proving the stated
complex version on the common domain.
\end{proof}

\section{Quadratic feedback: exact blocks and uniform remainders}
\label{sec:quadratic}

\subsection{A normalized expansion with nonnegative frequencies}
Take $\lambda_j=j^2$ and $P=\ee^uQ$.  In an atom from
\eqref{eq:frequency}, the new frequency is
\begin{equation}
 a_{k,\mathbf m}=b_{k,\mathbf m}+1
    =\sum_{i=1}^\ell m_i(m_i+1).
 \label{eq:quadratic-frequency}
\end{equation}
Because the parts sum to $k-1$,
\begin{equation}
 0\le a_{k,\mathbf m}\le k(k-1).
 \label{eq:quadratic-support}
\end{equation}
The upper bound follows from $\sum m_i^2\le(\sum m_i)^2$.
Thus there are explicitly computable polynomials $H_k(v)$ with
\begin{equation}
 P(u)=\sum_{k\ge1}u^kH_k(\ee^u),
 \qquad \sum_a|[v^a]H_k(v)|\le\sh_k,
 \qquad \deg H_k\le k(k-1).
 \label{eq:quadratic-block}
\end{equation}
The first four are
\begin{align*}
 H_1(v)&=1, & H_2(v)&=-v^2,\\
 H_3(v)&=2v^4-v^6, & H_4(v)&=-5v^6+5v^8-v^{12}.
\end{align*}
They are not the ordinary Taylor coefficients $a_n$ of $P$; the two
levels of expansion must be kept distinct.

Since all frequencies in \eqref{eq:quadratic-block} are nonnegative,
this block expansion converges uniformly on every closed left half-disc
$|u|\le r_0<\rho$.  A directly usable tail bound, for $r<1/6$, is
\begin{equation}
 \left|P(u)-\sum_{k=1}^K u^kH_k(\ee^u)\right|
 \le\frac{(6r)^{K+1}}{4(1-6r)},
 \qquad \Ree u\le0,\quad |u|\le r.
 \label{eq:geometric-tail}
\end{equation}
For $Q$, multiply this bound by $\ee^r$.
This is a convergent block approximation to a function with a divergent
ordinary Taylor expansion.  There is no contradiction: each block
retains an entire exponential instead of Taylor expanding it indefinitely.

\subsection{Proof of the uniform remainder certificate}
For $m\ge1$, the integral form of Taylor's remainder says
\[
 \ee^z-\sum_{j=0}^{m-1}\frac{z^j}{j!}
  =\frac{z^m}{(m-1)!}\int_0^1(1-t)^{m-1}\ee^{tz}\,\dd t.
\]
If $\Ree z\le0$, its modulus is at most $|z|^m/m!$.
Apply this to each exponential in block $k<N$, using $m=N-k$ and
\eqref{eq:quadratic-support}.  The total error from this block is at most
\[
 r^N\sh_k\frac{[k(k-1)]^{N-k}}{(N-k)!}.
\]
The unexpanded blocks $k\ge N$ contribute at most
\[
 \sum_{k\ge N}\rho^{-k}r^k
 \le r^N\frac{\rho^{-N}}{1-r_0/\rho}.
\]
The finite Taylor sums assemble to exactly
$\sum_{n<N}a_nu^n$, because the coefficient at degree $n$ uses only
$k\le n$.  Summing the bounds proves \eqref{eq:remainder} and
\eqref{eq:MN}.  Notice that this estimate is uniform all the way to
the imaginary boundary of the half-disc, where ordinary exponential
damping disappears.

\section{Asymptotics of the majorant and the deeper flat scale}
\label{sec:majorant}

\subsection{A discrete saddle for the positive majorant}
We now prove \eqref{eq:MN-asymp}, including its linear constant.
For $2\le k<N$ define
\[
 T_{N,k}=\sh_k\frac{[k(k-1)]^{N-k}}{(N-k)!}.
\]
The $k=1$ summand vanishes.  The logarithm of a sum of at most $N$ positive
terms differs from its largest logarithm by at most $\log N$.
The geometric-tail term in \eqref{eq:MN} has logarithm only $O(N)$.
It will be smaller than the saddle contribution.

Put $t=k/N$ and $c=\rho^{-1}$.  The bounds
$1\le\sh_k\le c^k$, $k(k-1)\le k^2$, and the uniform elementary
Stirling estimate give
\begin{equation}
 \frac1N\log T_{N,k}
 \le(1-t)\bigl(\log N+2\log t-\log(1-t)+1\bigr)
       +t\log c+O(\log N/N).
 \label{eq:saddle-upper}
\end{equation}
The error bound is uniform for $1\le N-k\le N$.
Set $a=t\log N$.  Subtracting
$\log N-2\log\log N$ from the right side of
\eqref{eq:saddle-upper} gives
\begin{equation}
 2\log a+1-a+R(t)+O(\log N/N),
 \label{eq:saddle-shape}
\end{equation}
where
\[
 R(t)=-2t\log t-(1-t)\log(1-t)-t+t\log c.
\]
The function $R$ is bounded on $[0,1]$, with its endpoint values interpreted
continuously, and $R(t)\to0$ as $t\to0$.

The function $2\log a+1-a$ tends to minus infinity both as
$a\downarrow0$ and as $a\to\infty$.  Consequently any indices capable of
maximizing \eqref{eq:saddle-shape} must have $a$ in a fixed compact
subinterval of $(0,\infty)$.  On such a set,
$t=a/\log N\to0$ uniformly, and hence $R(t)=o(1)$.
The limiting optimization is therefore exact:
\[
 \max_{a>0}(2\log a+1-a)=2\log2-1=\log4-1,
 \qquad a=2.
\]
For a matching lower bound choose
$k=\lfloor2N/\log N\rfloor$.  Use $\sh_k\ge1$ and Stirling's formula.
Replacing $k(k-1)$ by $k^2$ changes the logarithm by
$(N-k)\log(1-1/k)=O(\log N)=o(N)$.
The same computation without $t\log c$ therefore attains
$\log4-1+o(1)$.  The maximum and then the positive sum have precisely the
asymptotic \eqref{eq:MN-asymp}.  This completes the proof of
Theorem~\ref{thm:remainder}.

\begin{corollary}[Uniform Gevrey--1 of zero type]
For every $H>0$ and every fixed $r_0<\rho$ there is a constant $C_H$ such
that the remainder in \eqref{eq:remainder} is at most
$C_H H^N N! |u|^N$ for all $N\ge2$ and all $u$ in the stated half-disc.
\end{corollary}
\begin{proof}
From \eqref{eq:MN-asymp} and Stirling,
\[
 \frac1N\log\frac{M_N}{N!}
 =-2\log\log N+\log4+o(1)\longrightarrow-\infty.
\]
Thus $M_N\le H^N N!$ for sufficiently large $N$; enlarge $C_H$ for the
finitely many remaining indices.
\end{proof}

\subsection{Optimizing the proved bound}
\begin{proof}[Proof of Theorem~\ref{thm:optimal}]
Let $L=\log(1/r)$ and $N=\lfloor L^2/(4r)\rfloor$.  Then
\[
 \log N=L+2\log L-\log4+o(1),
 \qquad \log\log N=\log L+o(1).
\]
Using \eqref{eq:MN-asymp},
\begin{align*}
 \log(r^N M_N)
 &=N\bigl(-L+\log N-2\log\log N+\log4-1+o(1)\bigr)\\
 &=-\frac{L^2}{4r}(1+o(1)).
\end{align*}
The finite bound \eqref{eq:remainder} is uniform over the half-circle,
so this proves \eqref{eq:optimal}.  Multiplication by $\ee^r$ contributes
only $r$ to the logarithm of the error bound, which is absorbed by an
arbitrarily small loss in $\varepsilon$.
\end{proof}

Under the transseries substitution $u=-\ee^{-X}$, $X\to+\infty$, the
remainder scale becomes
\begin{equation}
 \exp\!\left[-\left(\frac14-o(1)\right)X^2\ee^X\right].
 \label{eq:deep-flat}
\end{equation}
This is invisible to every fixed exponential monomial $\ee^{-nX}$.
It is an analytically proved error scale arising from an infinite
reorganization of a one-generator formal exponential transseries.
The conclusion is an upper bound; it does not by itself specify a new
nonzero transseries sector or prove that this scale appears with a
nonvanishing coefficient in a completed expansion.

\subsection{What the constant does and does not say}
There are three different quantities here: the growth of the actual
Taylor coefficients, the growth of the positive majorant $M_N$, and the
smallest actual error among all truncations.  The saddle calculation
proves a sharp asymptotic for the second quantity.  It proves an upper
bound for the third, but no matching lower bound.  The signed block
coefficients can cancel, so replacing ``majorant asymptotic'' by
``actual optimal remainder asymptotic'' would be unjustified.

Likewise, an entire Borel transform need not have acceptable growth in
every direction.  The earlier feedback draft establishes, for the forward
quadratic model, an entire ordinary Borel transform with excessive growth
on the positive ray \cite{feedback}.  Our negative-ray recovery does not
contradict that obstruction.  It supplies the missing favorable direction,
and it does not assert an angular neighborhood of favorable directions.

For finite computation, use \eqref{eq:MN} rather than the limiting
constant.  The convergence in \eqref{eq:MN-asymp} is slow: the packaged
floating-point diagnostics still place the centered value
$(\log M_N-N\log N+2N\log\log N)/N$ near $1.35$ at $N=10^5$, whereas its
proved limit is $\log4-1\approx0.38629$.  This is a reason to retain
explicit certificates, not a discrepancy with a limit containing an
unspecified $o(N)$ term.

\section{Examples, exact computations, and inverse certificates}
\label{sec:verification}

\subsection{Examples across the threshold}
For $\lambda_j=j^2(\log(j+1))^\beta$, the conclusion is particularly simple:
\begin{center}
\begin{tabular}{@{}lll@{}}
\toprule
Slope regime & Fine negative summation & Guaranteed inverse Borel domain\\
\midrule
$\beta<2$ & Yes & Entire plane\\
$\beta=2$ & Yes & $|\zeta|+\Ree\zeta<1/2$\\
$\beta>2$ & No & No holomorphic Borel germ at zero\\
\bottomrule
\end{tabular}
\end{center}
The failure in the last row holds for both $\U$ and $\Q$.  Nevertheless,
Propositions~\ref{prop:all-slopes} and \ref{prop:ray} still construct the
literal negative-side inverse and its continuous ray representation.
The criterion is a limsup criterion: slopes exceeding the threshold by
unbounded factors on an arbitrarily sparse subsequence already prevent
fine Borel summability.  Eventual monotonicity or regular variation of the
slopes is not assumed.

For comparison, the convergence threshold is much smaller:
$\U$ and $\Q$ are convergent at zero exactly when $\lambda_j=O(j)$,
as already recorded in \cite{feedback}.  Here is a short independent
check.  If $\lambda_j\le Lj$, the kernel is holomorphic for
$|q|\ee^{L|u|}<1$ and the analytic implicit-function theorem applies at
$(0,0)$.  Conversely, if $u_n\le CH^n$, use \eqref{eq:lower} with $m=j$
to obtain $\lambda_j\le C^{1/j}H^2(j!)^{1/j}=O(j)$.
Convergent reversion preserves analyticity.  Thus in the quadratic model
both original Taylor series diverge, even though the inverse Borel
transform constructed here is entire.

Two soluble cases test the normalizations.  With all slopes zero,
$U(q)=q/(1-q)$ and $Q(u)=u/(1+u)$.  With $\lambda_j=j$,
\begin{equation}
 Q(u)=\frac{u\ee^{-u}}{1+u}.
 \label{eq:linear-check}
\end{equation}
The verification program checks these against a separate forward
coefficient recurrence.

\subsection{Independent exact coefficient checks}
The accompanying program computes the forward coefficients from
$E_j(q)=\exp(j^2U(q))$ using
\begin{equation}
 [q^m]E_j=\frac{j^2}{m}
             \sum_{i=1}^m i u_i[q^{m-i}]E_j,
 \qquad u_n=\sum_{j=1}^n[q^{n-j}]E_j.
 \label{eq:forward-rec}
\end{equation}
The first formula follows by differentiating $E_j$; the second is the
feedback equation.  This computation is independent of the inverse-block
construction.  The latter uses dynamic programming over compositions,
retaining exact rational coefficients and integer frequencies.  The first
ordinary coefficients are as follows.
\begin{center}
\begin{tabular}{@{}r r r r r@{}}
\toprule
$n$ & $u_n$ in $\U$ & $a_n$ in $\Pp$ & $v_n$ in $\Q$ & $\sh_n$\\
\midrule
1 & $1$ & $1$ & $1$ & 1\\
2 & $2$ & $-1$ & $-2$ & 1\\
3 & $15/2$ & $-1$ & $1/2$ & 3\\
4 & $107/3$ & $-1$ & $-2/3$ & 11\\
5 & $1555/8$ & $-13/3$ & $-29/8$ & 45\\
6 & $17362/15$ & $-49/3$ & $-743/60$ & 197\\
7 & $5288311/720$ & $-1099/15$ & $-8491/144$ & 903\\
8 & $15414767/315$ & $-16207/45$ & $-92726/315$ & 4279\\
\bottomrule
\end{tabular}
\end{center}
The executed exact-arithmetic checks verify both
$\U\circ\Q=u$ and $\Q\circ\U=q$ through degree 24, the recurrence of
$S$, the block-frequency bounds, and the uncancelled-mass bounds.  The
zero- and linear-slope tests also pass through degree 24.  These are finite
checks of formulas and implementation, not substitutes for the preceding
proofs.

\subsection{Three rational enclosures of the inverse function}
The block series gives numerical certificates without first estimating a
small divergent-series remainder.  For rational $r>0$, enclose
$t=\ee^{-r}$ between consecutive Taylor polynomials.  Taylor's integral
remainder gives the direction of the two inequalities exactly.
Evaluate the finite polynomial
$\sum_{k=1}^{K}(-r)^kH_k(t)$ by rational interval arithmetic, append the
rational tail bound \eqref{eq:geometric-tail}, and divide by the positive
interval for $t$.  This encloses $Q(-r)=P(-r)/t$.

For $K=14$, the executed program used Taylor degrees 36 and 37 for $t$,
then rounded the final rational endpoints outward to 24 decimal places.
Every displayed endpoint below is therefore an exact terminating rational,
not a rounded floating-point claim.
\begin{center}\small
\begin{tabular}{@{}c l l@{}}
\toprule
$r$ & Lower bound for $Q(-r)$ & Upper bound for $Q(-r)$\\
\midrule
$1/20$ & $-0.055065695744769945932257$ & $-0.055065684970061950441109$\\
$1/40$ & $-0.026258040219395119101712$ & $-0.026258040219131013674068$\\
$1/100$ & $-0.010200506315988279498804$ & $-0.010200506315988279246191$\\
\bottomrule
\end{tabular}
\end{center}
The corresponding widths are less than $1.078\times10^{-8}$,
$2.642\times10^{-13}$, and $2.527\times10^{-19}$, respectively; the
unrounded exact widths and all certificate inputs are in the data files.

\subsection{A certified derivative bound and error transport}
\begin{proposition}[A concrete inverse-error certificate]
\label{prop:errortransport}
For quadratic feedback and $u\in[-1/20,0]$, the inverse branch obeys
(the derivative at zero is interpreted from the left)
\begin{equation}
 |Q(u)|<1/10,\qquad Q'(u)\ge\frac{619}{900}>0.
 \label{eq:derivative}
\end{equation}
If $u_*$ and $v$ lie in this interval and $Q(u_*)=q_*$, then
\begin{equation}
 |u_*-v|\le\frac{900}{619}|Q(v)-q_*|.
 \label{eq:transport}
\end{equation}
If $Q(v)$ is known only through an approximation $\widetilde Q(v)$ with
certified error $\eta$, its residual in \eqref{eq:transport} can be replaced
by $|\widetilde Q(v)-q_*|+\eta$.  Existence of $u_*$ must first be checked
by a bracket; an error bound alone is not an existence proof.
\end{proposition}
\begin{proof}
For $r\le1/20$, \eqref{eq:C-bounds} gives
\[
 |Q(u)|\le\ee^r S(r)
 \le\frac1{1-r}\left(r+\frac{9r^2}{1-6r}\right)<\frac1{10}.
\]
The same estimate holds on the closed left half-disc.  It guarantees
that \eqref{eq:literal}, initially proved near zero, continues throughout
its open half-disc by the identity theorem, since the $q$-series stays
inside $|q|<1$.

Let $R=1/10$.  On the indicated real interval the differentiated kernel
satisfies
\[
 |\Phi_u(Q(u),u)|\le\sum_{j\ge1}j^2R^j
   =\frac{R(1+R)}{(1-R)^3}=\frac{110}{729},
\]
and
\[
 \Phi_q(Q(u),u)\ge\ee^u-\sum_{j\ge2}jR^{j-1}
 \ge\frac{19}{20}-\frac{19}{81}>0,
 \qquad \Phi_q(Q(u),u)\le\frac{100}{81}.
\]
These derivatives are justified by absolute convergence of their
geometric majorants.  Differentiating $u=\Phi(Q(u),u)$ yields
\[
 Q'(u)=\frac{1-\Phi_u}{\Phi_q}
 \ge\left(1-\frac{110}{729}\right)\frac{81}{100}
 =\frac{619}{900}.
\]
The mean-value theorem proves \eqref{eq:transport}; the approximate
residual follows from the triangle inequality.
\end{proof}

As an executed application, take $q_*=-1/40$.  The first rational
enclosure above and $Q(0)=0$ bracket this target, so a unique
$u_*=U(-1/40)$ lies in $[-1/20,0]$.  Choose the rational center
$v=-0.02385490927750544$, without assuming its accuracy.  Enclosing $Q(v)$
by 14 blocks and Taylor degrees 12 and 13 for $\ee^v$, then applying
\eqref{eq:transport} and outward rounding, gives the exact certificate
\begin{equation}
 \boxed{\quad
 -0.023854909277599602576607
 \ \le U(-1/40)\le\
 -0.023854909277411277423393.
 \quad}
 \label{eq:U-cert}
\end{equation}
The width is $1.88325153214\times10^{-13}$.
The center was selected numerically, but its certificate uses only rational
arithmetic and Proposition~\ref{prop:errortransport}.

\subsection{Floating-point checks are separately labeled}
At 70 decimal digits, the program compares the 24-block inverse
approximation with a numerical root of a 179-term literal feedback sum
at $r=0.01,0.025,0.05$.  It also numerically integrates four independent
Bessel kernels and compares the integrals with
\eqref{eq:kernel-laplace}.  All four kernel comparisons have absolute
discrepancy below $10^{-60}$ in that run.  These calculations are
floating-point diagnostics, not interval proofs of the roots or of the
integrals.  The rational enclosures and residual-transport certificate
are separate computations and do not depend on those comparisons.

The majorant diagnostic program evaluates a positive sum in log-space
and records the slow movement of its maximizing block.  It uses the
Schr\"oder recurrence
\[
 k\sh_k=3(2k-3)\sh_{k-1}-(k-3)\sh_{k-2},\qquad k\ge3,
\]
in ratio form.  Its output is explicitly marked non-certified; the proof
of the limiting saddle does not use that output.

\section{Further research questions and proposed directions}
\label{sec:questions}

The following questions are not claimed solved in this article.  They
are framed so that a successful answer would add information not already
contained in the fine summation theorem.

\subsection{Angular summability of the quadratic model}
Determine the set of angles $\theta$ for which the entire functions
$\B\Q(te^{i\theta})$ and $\B\U(te^{i\theta})$ have at most exponential
growth as $t\to\infty$, with the uniform neighborhood needed for angular
summability.  In particular, does any open arc containing $\pi$ work?
The present fixed-width tube does not answer that question.  The
uncancelled Bessel estimate deteriorates off the negative ray, so a proof
of an open arc would have to exploit cancellation or a different exact
representation.  A proof that no such arc exists would identify a genuine
boundary between fine and angular summability in a concrete nonlinear
implicit model.
% ed. (2026-09-29): open-arc question answered negatively by the later natural-boundaries package.
\begin{ednote}
Answered negatively for $\lambda_j=j^2$ by the later package
\nolinkurl{Natural_Boundaries_Quadratic_Exponential_Feedback/} (its
Theorem~\texttt{thm:main}), which cites this subsection. The literal inverse $Q$ on
$|u|\le1/32$, $\Ree u\le0$ is smooth up to the imaginary diameter, every
point $it$ with $|t|<1/32$ is a natural-boundary point, and $\B\Q$ is
entire and exponentially bounded on every fixed-width tube around the
negative ray (re-proving the fine result here in this case) but in no
wedge around it; hence neither $\Q$ nor $\U$ is angularly $1$-summable in
direction $\pi$. For rational amplitude generating functions with
eventually quadratic slopes, angular summability holds exactly when the
action support is finite (its Theorem~\texttt{thm:rational-main}). The set of
admissible individual rays and the growth in shrinking sectors remain
open. That package has not been independently reviewed.
% ed. (2026-09-29): added on filing batch 50; a second, independent negative answer.
It is also answered negatively, independently, by the later package
\nolinkurl{Natural_Boundaries_Survive_Nonlinear_Feedback/} (batch~50; its
Theorems~\texttt{thm:main} and \texttt{thm:quadratic-main}), which proves the
natural boundary for $\lambda_j=j^d$ and every integer $d\ge2$; not
reviewed.
\end{ednote}

\subsection{The actual Borel boundary at critical slope growth}
For $\lambda_j\sim a(j\log j)^2$, $a>0$, determine the singularities of
$\B\Q$ and whether any part of the boundary
$|\zeta|+\Ree\zeta=1/(2a)$ is an actual obstruction.  Our theorem proves
normal convergence strictly inside; it does not prove failure outside.
A related question is the exact factorial-normalized type of $\Q$ and its
relation to the forward type.  Class membership is invariant under
reversion, but the elementary proof here does not establish a numerical
type identity.  This is a useful place for a more detailed literature
comparison before asserting novelty of any type-preservation theorem.
% ed. (2026-09-29): type part answered by the later weighted-type package; boundary part only at the vertex (editorial deduction).
\begin{ednote}
The type part is answered by the later package
\nolinkurl{Sharp_Weighted_Type_Formal_Reversion/}, which cites this
package's README: $T_s(\Q)=T_s(\U)$ for every $s>0$ (its
Theorem~\texttt{thm:inversefeedback}, using the regularity article's forward type
identity as an imported input), and the radius of convergence of $\B\Q$
is exactly $1/(4A)$ (its Corollary~\texttt{cor:borelradius}). Since every point of
the circle $|\zeta|=1/(4A)$ other than $\zeta=1/(4A)$ lies in the region
\eqref{eq:parabola}, the two results together imply, for $0<A<\infty$,
that the vertex $\zeta=1/(4A)$ is a singular point of $\B\Q$. This
deduction is editorial and is stated in neither package; the rest of the
boundary is not decided.
\end{ednote}

\subsection{Signed coefficient asymptotics and actual least error}
Find an asymptotic equivalent, including signs and prefactors, for the
ordinary coefficients of $\Pp$ or $\Q$ in the quadratic case.  The
positive majorant deliberately ignores cancellation between compositions
with different lengths.  A signed saddle analysis may reveal a smaller
coefficient scale or a different correction to the dominant block.
The stronger analytic target is a matching lower bound for the best
Taylor-truncation error, or a full description of its exceptional zeros.
Only then would it be legitimate to interpret a constant such as $1/4$
as the actual least-error action rather than a certificate constant.
% ed. (2026-09-29): coefficient part claimed by two batch-49 packages (unreviewed); least-error part open.
\begin{ednote}
The coefficient part is claimed, without citing this subsection, by two
independent later packages that have not been reviewed:
\nolinkurl{Quadratic_Exponential_Feedback_After_Reversion/} and
\nolinkurl{Signed_Quadratic_Feedback_Inversion/} (Theorem~\texttt{thm:main} in
each) prove, for $\lambda_j=aj^2$, the signed equivalent
$v_n\sim-S_n\exp(-(1+1/a)r_n)$ (with changed constants under finitely
many exceptional weights and slopes) conjectured in
\nolinkurl{Finite_Core_Universality_Exponential_Feedback/}
(\texttt{conj:quadratic-inverse}), with that package's $r_n$ and $S_n$. A matching
lower bound for the best truncation error is not claimed by any package.
% ed. (2026-09-29): added on filing batch 50; a third claimed proof and the least formal term.
A third independent claimed proof, under the hypotheses of the first, is
\nolinkurl{Signed_Condensation_Sharp_Quadratic_Reversion/} (batch~50), which
also gives the least formal inverse term (its Theorem~\texttt{thm:least}):
at $a=1$ its exponent $-\log^2(1/t)/(4t)$ has the constant $1/4$ of
Theorem~\ref{thm:optimal}. That is a statement about formal terms and gives
no analytic lower bound.
\end{ednote}

\subsection{Other powers and generalized summation kernels}
For regularly varying slopes $\lambda_j\sim a j^p$ with $p>1$, construct
analogues of the inverse kernel formula for moment Borel transforms with
weights $\Gamma(1+(p-1)n)$.  Determine which directions have sufficient
kernel damping, and whether the literal left-side inverse is recovered.
For $p>2$ the ordinary Borel germ need not exist, so changing the summation
order is essential rather than cosmetic.  The quadratic proof suggests
that keeping exponential blocks intact may remain useful, but it does
not establish a generalized summability theorem for those other orders.

\subsection{Complex slopes, signs, and loss of positivity}
Replace the nonnegative slopes by complex slopes, perhaps contained in a
fixed sector, and identify sufficient geometric conditions for an
oscillatory or damped Borel direction.  Then seek necessity statements.
Our forward lower bound uses positivity and generally fails when
coefficients cancel; the negative-ray Bessel bound also uses the slopes'
common real direction.  Both parts of the equivalence must therefore be
revisited.  A common half-plane of damping for the literal kernel may
suffice for analytic inversion without sufficing for fine summability.

\subsection{Weighted actions and joint amplitude--slope thresholds}
Study
\[
 U(q)=\sum_{j\ge1}c_jq^j\ee^{\lambda_jU(q)},\qquad c_1=1,
\]
with positive amplitudes that can decay rapidly.  Determine the sharp
joint condition on $(c_j,\lambda_j)$ for fine summability and for
convergence.  The inverse-block formula survives with additional products
of the $c_j$, but the mass generating function and the positive coefficient
lower bound both change.  Very small amplitudes can suppress isolated
large slopes, so a criterion involving only $\sup\lambda_j/j$ or a bare
slope limsup should not be expected in this setting.
% ed. (2026-09-29): answered for positive exponentially bounded amplitudes by a later package (batch 50).
\begin{ednote}
Answered for positive exponentially bounded amplitudes by the later
package \nolinkurl{Amplitude_Slope_Compensation_Feedback_Transseries/}
(batch~50; its Theorems~\texttt{thm:classification} and
\texttt{thm:fine}; not reviewed): with $c_1=1$ and $d_j=-\log c_j$, the
solution converges exactly when $\lambda_j=O(j+d_j)$, and for $s=1$ the
solution and its inverse are finely summable in direction $\pi$ exactly when
$\lambda_j^{1/2}=O(j\log(j+1)+d_j)$. At $c_j=1$ this is the equivalence
proved here.
\end{ednote}

\subsection{Canonical recovery beyond every finite Gevrey class}
For superpolynomial slopes, the all-slope ray construction still produces
a continuous, exponentially bounded transform recovering the literal
inverse.  Relate this transform to acceleration, moment summability, or
other summation procedures defined from formal coefficients.  The main
tasks are to specify a summation class and prove that its canonical value
coincides with the literal implicit branch.  Calling the continuous ray
function a Borel germ would evade, rather than solve, this problem.
Uniqueness should be proved within the chosen summation class and not
inferred merely from having a common Poincar\'e expansion.

\subsection{Sharper certified algorithms and automatic normalization}
Develop an algorithm that recognizes when a normalization such as
$P=\ee^uQ$ removes all negative block frequencies, then selects between
geometric block truncation and ordinary Taylor truncation using certified
cost and error estimates.  Improving the uncancelled mass bound could
substantially reduce the number of blocks needed in rational certificates.
A further goal is an interval-verified Borel--Laplace quadrature directly
from the Bessel representation, with rigorously bounded tails and a
verified equivalence to the simpler block evaluator.

\section{A formalization roadmap and proof obligations}
\label{sec:lean}

The repository's existing formal scales and support infrastructure are
relevant context, but they do not automatically formalize any analytic
summation assertion proved here.  A practical staged development would
separate finite algebra from measure, complex-analysis, and uniform-limit
arguments.

The first stage should formalize the composition-indexed coefficient
measures as finitely supported maps, prove the exact frequency formulas,
and establish the mass recurrence and generating-function identities.
The quadratic bound $0\le a\le k(k-1)$ and the two composition checks are
finite algebraic targets.  They require no Borel integration theory.
An exact certificate checker can verify the rational interval arithmetic
using proved exponential Taylor bounds and the geometric tail theorem.

The second stage should establish the normalized beta moments and
\eqref{eq:poisson}, followed by the one-block Laplace identity.  The proof
must distinguish an absolutely integrable Taylor series for one fixed
frequency from the much stronger absolute-integrability estimate used
when summing infinitely many blocks.  Confusing those two exchanges would
invalidate the all-slope construction.

The third stage should formalize the compact-normal convergence estimate
\eqref{eq:master-bound} and the fixed-strip split at $k=s^2$.  Only after
those are in place should the implementation identify the formal Borel
coefficients of the limiting function.  The nonlinear inverse theorem
then belongs in a reusable library for half-strip Borel minors and
Laplace summation, rather than being hidden in a feedback-specific lemma.

Finally, the logarithmic saddle proof should have an explicit statement
about the \emph{positive majorant}, with a separate theorem transporting
it to an upper bound on a function error.  This distinction is valuable
for formalization: it prevents a later user from inadvertently rewriting
an inequality as an asymptotic equivalence for an oscillating remainder.
No Lean files are supplied with this article, and no claim is made that
these stages have been implemented or checked.

\section{Conclusion}
The countable feedback model provides a concrete example in which formal
support, coefficient growth, directional summation, and uniform error
control are genuinely different layers.  Inverting the dependent variable
first preserves enough exponential structure to bridge several of them.
The resulting Bessel representation makes the negative direction
accessible exactly up to the known Gevrey--1 slope threshold, identifies
the summed functions with the literal implicit branch, and yields a
stronger quantitative theorem in the quadratic case.

The proven progress is therefore specific: a negative-direction fine
summation theorem, a guaranteed inverse Borel continuation region, a
closed-half-disc logarithmic-Gevrey certificate, and reproducible exact
inverse bounds.  Angular summability, maximal continuation, global
resurgence, and actual optimal-remainder asymptotics remain separate
research questions.  These boundaries are part of the result, not
qualifications to be removed in a future summary.

\appendix
\section{Reproducibility and package contents}
\label{sec:reproduce}
The archive contains the self-contained source \texttt{article.tex}, the
compiled \texttt{article.pdf}, the two programs in \texttt{code/}, and
recorded data in \texttt{data/}.  The primary exact-arithmetic program uses
Python's standard-library fractions.  If \texttt{mpmath} is installed, it
also runs the separately labeled numerical diagnostics.  The majorant
program uses only the standard library.

The exact commands used to generate the data and document are
\begin{quote}\small\ttfamily
python code/verify.py --order 24\\
python code/majorant.py\\
pdflatex -interaction=nonstopmode -halt-on-error article.tex\\
pdflatex -interaction=nonstopmode -halt-on-error article.tex\\
pdflatex -interaction=nonstopmode -halt-on-error article.tex
\end{quote}
The recorded run used Python 3.13.5 and mpmath 1.3.0 at 70 digits for the
floating-point checks.  The scripts accept ordinary local execution and
make no network requests.  No source file or branch in ProveIt was modified.

The certificate files store rational numerators and denominators in
addition to decimal displays.  The two composition identities are checked
by truncated polynomial composition, not by comparing a few numerical
function values.  Re-running at a larger order extends those finite checks
but does not change the status of the analytic proofs.

\section{An explicit coefficient-extraction algorithm}
For the quadratic blocks, let $D_{s,\ell}(a)$ count ordered compositions
of $s$ into $\ell$ positive parts whose frequency
$\sum m_i(m_i+1)$ equals $a$.  Start with $D_{0,0}(0)=1$, and use
\begin{equation}
 D_{s,\ell}(a)=\sum_{m=1}^{s}
                     D_{s-m,\ell-1}(a-m(m+1)).
 \label{eq:dp}
\end{equation}
Absent indices have value zero.  Then
\begin{equation}
 [v^a]H_k(v)=\frac1k\sum_{\ell=1}^{k-1}
          (-1)^\ell\binom{k+\ell-1}{\ell}D_{k-1,\ell}(a),
 \qquad k\ge2.
 \label{eq:dp-H}
\end{equation}
Finally,
\begin{equation}
 a_n=\sum_{k=1}^n\sum_a[v^a]H_k(v)\frac{a^{n-k}}{(n-k)!},
 \qquad
 v_n=\sum_{j=0}^{n-1}\frac{(-1)^j}{j!}a_{n-j}.
 \label{eq:coeff-algorithm}
\end{equation}
The convention $0^0=1$ handles $H_1=1$.  All these expressions are finite
and rational.  Equation~\eqref{eq:dp} aggregates compositions with the same
frequency, avoiding explicit enumeration of all ordered compositions.
It is a reproducible reference algorithm, not a claim of asymptotically
optimal computational complexity.

\clearpage
\begin{thebibliography}{9}
\bibitem{project}
Vladimir Reshetnikov, \emph{ProveIt}, repository snapshot\newline
\texttt{a795fcffa4b3ece22761d81fbed3c43be8b81795}.\newline
\path{Analysis/Transseries/README.md} and\newline
\path{Analysis/Transseries/docs/series-and-transseries/README.md}.
Accessed 29 September 2026.
\href{https://github.com/VladimirReshetnikov/ProveIt/tree/a795fcffa4b3ece22761d81fbed3c43be8b81795/Analysis/Transseries}{Pinned repository directory}.

\bibitem{feedback}
\emph{A sharp regularity classification for countable exponential-feedback
transseries}, research draft prepared for Vladimir Reshetnikov,
29 September 2026.  Article and README in
\texttt{Exponential\_Feedback\_Regularity\_Classification/},
under the series-and-transseries directory of \cite{project}.
\href{https://github.com/VladimirReshetnikov/ProveIt/tree/a795fcffa4b3ece22761d81fbed3c43be8b81795/Analysis/Transseries/docs/series-and-transseries/Exponential_Feedback_Regularity_Classification}{Pinned research package}.
Cited as a repository research draft, not as a peer-reviewed or
machine-verified publication.

\bibitem{gessel}
Ira M. Gessel, \emph{Lagrange inversion}, Journal of Combinatorial Theory,
Series A \textbf{144} (2016), 212--249.
\href{https://doi.org/10.1016/j.jcta.2016.06.018}{doi:10.1016/j.jcta.2016.06.018}.
\href{https://arxiv.org/abs/1609.05988}{arXiv:1609.05988}.

\bibitem{dlmf}
NIST Digital Library of Mathematical Functions,
\emph{\S10.9, Integral Representations}, especially the Poisson integral
10.9.4.  Accessed 29 September 2026.
\url{https://dlmf.nist.gov/10.9}.

\bibitem{sauzin}
David Sauzin, \emph{Introduction to 1-summability and resurgence},
2014, 127 pp.  In particular, Sections 7 and 9 distinguish fine
half-strip summability from summability in an arc; Sections 13 and 17
explain nonlinear operations and inversion.
\href{https://arxiv.org/abs/1405.0356}{arXiv:1405.0356}.

\end{thebibliography}
\end{document}
